\documentclass[runningheads]{llncs}
\usepackage{graphicx}
\usepackage{booktabs}
\usepackage{longtable}
\usepackage{float}
\usepackage{array}
\usepackage{microtype}
\usepackage{amsmath}
\usepackage[section]{placeins}
\usepackage[T5,T1]{fontenc}
\usepackage[hidelinks]{hyperref}
\emergencystretch=1em
\newcommand{\method}{CEARF-N}
\newcommand{\rr}{\rho}
% Generated by summarize_dynamic_beta.py; do not edit manually.
\newcommand{\DynamicSeedCount}{3}
\newcommand{\DynamicVideoRtwenty}{.14735}
\newcommand{\DynamicVideoBetaMean}{.572}
\newcommand{\DynamicVideoBetaWithinSD}{.013}
\newcommand{\DynamicVideoNetVsGlobal}{15}
\newcommand{\DynamicVideoLowBetaAdv}{-.01026}
\newcommand{\DynamicVideoHighBetaAdv}{.00589}
\newcommand{\DynamicVideoAllocationGap}{.01614}
\newcommand{\DynamicVideoOOFRows}{4000 (0.0)}
\newcommand{\DynamicVideoActionable}{1666 (25.4)}
\newcommand{\DynamicVideoActionableShare}{.4165 (.0063)}
\newcommand{\DynamicVideoGlobalBeta}{.57748 (.00705)}
\newcommand{\DynamicVideoCandidatePool}{94,762}
\newcommand{\DynamicVideoDeclaredValidation}{5,000}
\newcommand{\DynamicVideoTargetsRemoved}{0}
\newcommand{\DynamicVideoDeclaredOOFOverlap}{0}
\newcommand{\DynamicVideoCandidateOOFOverlap}{5,000}
\newcommand{\DynamicVideoEligible}{89,762}
\newcommand{\DynamicVideoProfileHash}{623fdf1e}
\newcommand{\DynamicVideoGateHash}{3df29839}
\newcommand{\DynamicVideoRescues}{4451 (29.9)}
\newcommand{\DynamicVideoDamage}{1766 (43.3)}
\newcommand{\DynamicVideoNetRescues}{2685 (19.1)}
\newcommand{\DynamicVideoPrimaryU}{.11281 (.00021)}
\newcommand{\DynamicVideoContextMLPU}{.11303 (.00026)}
\newcommand{\DynamicVideoFullLinearU}{.11295 (.00032)}
\newcommand{\DynamicVideoFullMLPU}{.11310 (.00044)}
\newcommand{\DynamicVideoNoCrossExpertU}{.11297 (.00034)}
\newcommand{\DynamicVideoNoMemoryCertaintyU}{.11309 (.00026)}
\newcommand{\DynamicVideoKTenU}{.11219 (.00018)}
\newcommand{\DynamicVideoKTwentyU}{.11281 (.00021)}
\newcommand{\DynamicVideoKSixtyU}{.11163 (.00026)}
\newcommand{\DynamicBabyRtwenty}{.05669}
\newcommand{\DynamicBabyBetaMean}{.541}
\newcommand{\DynamicBabyBetaWithinSD}{.005}
\newcommand{\DynamicBabyNetVsGlobal}{-3}
\newcommand{\DynamicBabyLowBetaAdv}{.00245}
\newcommand{\DynamicBabyHighBetaAdv}{.00118}
\newcommand{\DynamicBabyAllocationGap}{-.00127}
\newcommand{\DynamicBabyOOFRows}{4000 (0.0)}
\newcommand{\DynamicBabyActionable}{845 (13.1)}
\newcommand{\DynamicBabyActionableShare}{.2113 (.0033)}
\newcommand{\DynamicBabyGlobalBeta}{.54335 (.00923)}
\newcommand{\DynamicBabyCandidatePool}{150,777}
\newcommand{\DynamicBabyDeclaredValidation}{5,000}
\newcommand{\DynamicBabyTargetsRemoved}{0}
\newcommand{\DynamicBabyDeclaredOOFOverlap}{0}
\newcommand{\DynamicBabyCandidateOOFOverlap}{5,000}
\newcommand{\DynamicBabyEligible}{145,777}
\newcommand{\DynamicBabyProfileHash}{b0fb3c22}
\newcommand{\DynamicBabyGateHash}{3556eb83}
\newcommand{\DynamicBabyRescues}{3889 (111.5)}
\newcommand{\DynamicBabyDamage}{1191 (55.9)}
\newcommand{\DynamicBabyNetRescues}{2698 (55.8)}
\newcommand{\DynamicBabyPrimaryU}{.04343 (.00032)}
\newcommand{\DynamicBabyContextMLPU}{.04345 (.00030)}
\newcommand{\DynamicBabyFullLinearU}{.04349 (.00041)}
\newcommand{\DynamicBabyFullMLPU}{.04348 (.00039)}
\newcommand{\DynamicBabyNoCrossExpertU}{.04345 (.00035)}
\newcommand{\DynamicBabyNoMemoryCertaintyU}{.04347 (.00041)}
\newcommand{\DynamicBabyKTenU}{.04261 (.00034)}
\newcommand{\DynamicBabyKTwentyU}{.04343 (.00032)}
\newcommand{\DynamicBabyKSixtyU}{.04354 (.00027)}
\newcommand{\DynamicDigiRtwenty}{.51701}
\newcommand{\DynamicDigiBetaMean}{.420}
\newcommand{\DynamicDigiBetaWithinSD}{.010}
\newcommand{\DynamicDigiNetVsGlobal}{-1}
\newcommand{\DynamicDigiLowBetaAdv}{-.01263}
\newcommand{\DynamicDigiHighBetaAdv}{-.01541}
\newcommand{\DynamicDigiAllocationGap}{-.00278}
\newcommand{\DynamicDigiOOFRows}{4000 (0.0)}
\newcommand{\DynamicDigiActionable}{3298 (8.2)}
\newcommand{\DynamicDigiActionableShare}{.8245 (.0020)}
\newcommand{\DynamicDigiGlobalBeta}{.41848 (.00284)}
\newcommand{\DynamicDigiCandidatePool}{133,724}
\newcommand{\DynamicDigiDeclaredValidation}{5,000}
\newcommand{\DynamicDigiTargetsRemoved}{5,000}
\newcommand{\DynamicDigiDeclaredOOFOverlap}{0}
\newcommand{\DynamicDigiCandidateOOFOverlap}{5,000}
\newcommand{\DynamicDigiEligible}{128,724}
\newcommand{\DynamicDigiProfileHash}{c0f021cd}
\newcommand{\DynamicDigiGateHash}{13a5b5af}
\newcommand{\DynamicDigiRescues}{3076 (55.8)}
\newcommand{\DynamicDigiDamage}{1693 (43.1)}
\newcommand{\DynamicDigiNetRescues}{1383 (86.6)}
\newcommand{\DynamicDigiPrimaryU}{.41332 (.00115)}
\newcommand{\DynamicDigiContextMLPU}{.41331 (.00108)}
\newcommand{\DynamicDigiFullLinearU}{.41345 (.00094)}
\newcommand{\DynamicDigiFullMLPU}{.41321 (.00105)}
\newcommand{\DynamicDigiNoCrossExpertU}{.41304 (.00093)}
\newcommand{\DynamicDigiNoMemoryCertaintyU}{.41357 (.00079)}
\newcommand{\DynamicDigiKTenU}{.41271 (.00081)}
\newcommand{\DynamicDigiKTwentyU}{.41332 (.00115)}
\newcommand{\DynamicDigiKSixtyU}{.41227 (.00074)}

% Generated by benchmark_dynamic_beta_inference.py; do not edit.
\newcommand{\DynamicRuntimeVideoQueries}{94,762}
\newcommand{\DynamicRuntimeVideoGateUs}{1.171}
\newcommand{\DynamicRuntimeVideoPostUs}{85.89}
\newcommand{\DynamicRuntimeVideoExpertMs}{3.710}
\newcommand{\DynamicRuntimeVideoTotalMs}{3.796}
\newcommand{\DynamicRuntimeBabyQueries}{150,777}
\newcommand{\DynamicRuntimeBabyGateUs}{1.184}
\newcommand{\DynamicRuntimeBabyPostUs}{87.57}
\newcommand{\DynamicRuntimeBabyExpertMs}{4.031}
\newcommand{\DynamicRuntimeBabyTotalMs}{4.119}
\newcommand{\DynamicRuntimeDigiQueries}{60,858}
\newcommand{\DynamicRuntimeDigiGateUs}{.774}
\newcommand{\DynamicRuntimeDigiPostUs}{80.60}
\newcommand{\DynamicRuntimeDigiExpertMs}{2.428}
\newcommand{\DynamicRuntimeDigiTotalMs}{2.509}
\newcommand{\DynamicRuntimeMaxBetaError}{$5.96\times10^{-8}$}
\newcommand{\DynamicRuntimeExactTopTwenty}{exact}


\begin{document}
\raggedbottom
\thispagestyle{empty}
\markboth{P. Q. Cuong et al.}{Supplementary Material for \method{}}
\begin{center}
{\Large\bfseries Supplementary Material for \method{}:\\
Bounded Out-of-Fit Dynamic Rank Allocation\\
for Sparse Next-Item Recommendation}

\vspace{0.6em}
Pham Quoc Cuong, Nguyen Ngoc Minh, Le Quang Minh,\\
and Nguyen Hoai Nguyet An

\vspace{0.35em}
{\small University of Information Technology, Ho Chi Minh City, Vietnam\\
Vietnam National University Ho Chi Minh City, Vietnam}
\end{center}

\vspace{0.6em}
\noindent This supplement specifies the complete \method{} protocol behind
the eight-page short paper: expert scoring, source-disjoint OOF construction,
continuous dynamic-$\beta$ optimization, controls, matched-query statistics,
and runtime verification. Neither declared-validation nor test labels fit or
select a fusion coefficient.

\section{Scope and Claim Ledger}
\begingroup\small
\begin{longtable}{p{.29\textwidth}p{.63\textwidth}}
\caption{What the paper does and does not claim.}\\
\toprule
Claim & Precise scope\\
\midrule
\endfirsthead
\multicolumn{2}{c}{\tablename\ \thetable{} -- continued}\\
\toprule
Claim & Precise scope\\
\midrule
\endhead
Training-only allocation &
The scalar prior and query-specific $\beta_q$ are fit only from targets of
source sessions removed before expert fitting. There is no $\beta$ grid.\\
Use of validation &
A stable hash first declares 5,000 validation queries from each loader's
candidate pool. This subset may audit the frozen system, and an earlier study
used validation to freeze the PASGR cell. It does not train, select, or
early-stop $\beta$.\\
Dynamic gain &
The primary comparison is dynamic versus an OOF-trained global or short/long
policy on identical expert ranks, not dynamic versus a weak endpoint. Video
has an interval-above-zero gain; Baby and Diginetica intervals span zero.\\
Assignment signal &
One fixed deterministic reassignment of the exact frozen $\beta_q$ multiset breaks
only query-to-beta matching. Video nDCG@20 and Diginetica nDCG@10 have
intervals above zero for this negative control; Baby is unresolved. Its
SHA-256 seed uses only the experiment seed and ordered query IDs, never
targets.\\
Metadata attribution &
PASGR uses a cache documented as E5-small~\cite{wang2022e5} on Video and TF--IDF/SVD teachers
on Baby and Diginetica. Comparisons with ID-only models are system-level;
matched-teacher Amazon results are labelled separately.\\
Expert sensitivity &
The NARM~\cite{li2017narm} swap changes only the neural expert inside the same allocator and is
reported as system-level sensitivity, not as metadata-matched architecture
attribution.\\
Generality &
Three domains support a bounded empirical finding, not a literature-wide
state-of-the-art or universal routing claim.\\
\bottomrule
\end{longtable}
\endgroup

\subsection{Methodological position}
The contribution is not a new encoder, RRF~\cite{cormack2009rrf}, or validation split. It is the
combination of frozen heterogeneous experts, source-disjoint OOF supervision,
target-free inference features, and a verifiable perturbation bound.
Out-of-fit meta-learning is rooted in stacked generalization
\cite{wolpert1992stacking}; our contribution is specifically the five-scalar,
bounded allocation law with a rank-only expert interface and its
frozen-expert protocol.
\begin{table}[h]
\caption{Allocation design space. ``OOF-safe'' means the allocator never
learns from predictions made on an expert's own fit sources.}
\centering\small
\begin{tabular}{lcccc}
\toprule
Family & Per-query & Frozen experts & OOF-safe & Bounded\\
\midrule
Global weighted rank fusion & No & Yes & N/A & N/A\\
In-fit stacking/router & Yes & Often & No & Usually no\\
OOF stacking & Yes & Yes & Yes & Usually no\\
Joint neural MoE & Yes & No & N/A & Usually no\\
\method{} & Yes & Yes & Yes & Yes\\
\bottomrule
\end{tabular}
\end{table}

\section{Question-to-Control Map}
This section maps each research question to the control that changes only one
axis of the proposed allocation protocol.
\begin{table}[h]
\caption{Each control changes one methodological axis while preserving the
declared train-only allocation boundary.}
\label{tab:question-control-map}
\centering\small
\begin{tabular}{
>{\raggedright\arraybackslash}p{.24\textwidth}
>{\raggedright\arraybackslash}p{.31\textwidth}
>{\raggedright\arraybackslash}p{.35\textwidth}}
\toprule
Question & Control & Interpretation\\
\midrule
Is fusion complementary? &
Memory-only, neural-only, equal mixing &
Separates evidence complementarity from allocator complexity.\\
Does per-query allocation matter? &
OOF global, short/long, head/tail, four buckets &
Compares one scalar, coarse policies, and continuous $\beta_q$.\\
Is the assignment itself useful? &
Fixed reassignment of primary $\beta_q$ &
Preserves the exact beta multiset while breaking query-to-beta matching.\\
Is the bound arbitrary? &
$\Delta\in\{.05,.10,.20\}$; no-prior/no-admission variants &
Tests modulation capacity and both shrinkage terms.\\
Which signals matter? &
Single-feature and drop-one gates; 16-feature linear/MLP &
Audits context, popularity/tail, and expert diagnostics.\\
Is the fusion implementation decisive? &
$k\in\{10,20,60\}$ and normalized CombSUM &
Tests the rank constant and a fixed-allocation score-space perturbation.\\
\bottomrule
\end{tabular}
\end{table}

\section{End-to-End Protocol}
\begin{enumerate}
\item Declare 5,000 validation queries per domain by stable hash before any
allocation fit; exclude their source identifiers from OOF construction.
\item From the remaining training sources, remove up to 5,000 complete
sessions: 1,000 for memory-profile locking and 4,000 for gate calibration.
\item Fit CEARF and PASGR on the inner sessions, then predict the two disjoint
OOF subsets. Profile selection uses only the first subset.
\item On actionable gate queries, learn the continuous
$\beta_{\mathrm{OOF}}$ from reciprocal-rank pairs, then fit the bounded
three-feature residual around that frozen prior.
\item Freeze the split fingerprints, expert identities, feature statistics,
scalar/gate states, and protocol flags before any declared-validation or test
prediction.
\item Within the allocation experiment, use declared validation only for
frozen-policy diagnostics. It cannot
change $\beta_{\mathrm{OOF}}$, the gate, $\Delta$, or the primary feature set.
\item Refit experts/indexes on the complete training data permitted for test,
retain the frozen allocator, rank the full catalogue, and save per-query ranks
for paired analysis.
\end{enumerate}

\section{Task, Data Construction, and Evaluation}
A query is a prefix $q=(i_1,\ldots,i_L)$ with one held-out next item $y_q$.
The catalogue is $\mathcal I$, and every reported ranking is over the full
catalogue. For Amazon~\cite{hou2024amazonreviews}, histories are ordered and split leave-last-out; items
already present in the query are masked. For Diginetica, the public session
split~\cite{diginetica2016} is retained and repeat consumption is allowed. Candidate generation may
operate at width 120 internally, but metrics inspect the final top 6, 10, or
20 positions. No sampled test negatives are used.

Each loader exposes a validation-candidate mapping before this experiment.
The study declares exactly 5,000 candidates per domain by stable hash before
any allocation fitting. On Diginetica, those 5,000 target events and incoming
transitions are removed from the tuning sessions; unselected candidate target
events remain ordinary training events. On Amazon, validation targets are
external to training, while source histories corresponding to unselected
candidates remain OOF-eligible. No unselected candidate is evaluated as part
of the declared validation subset.

\begin{table}[h]
\caption{Audited validation boundary. ``Candidate--OOF'' counts source
overlap with the loader's full candidate mapping; source histories outside the
declared 5,000 remain OOF-eligible. Declared--OOF is zero.}
\label{tab:validation-boundary}
\centering\small
\begin{tabular}{lrrrrr}
\toprule
Domain & Candidate pool & Declared & Removed & Declared--OOF &
Candidate--OOF\\
\midrule
Video Games & \DynamicVideoCandidatePool & \DynamicVideoDeclaredValidation &
\DynamicVideoTargetsRemoved & \DynamicVideoDeclaredOOFOverlap &
\DynamicVideoCandidateOOFOverlap\\
Baby Products & \DynamicBabyCandidatePool & \DynamicBabyDeclaredValidation &
\DynamicBabyTargetsRemoved & \DynamicBabyDeclaredOOFOverlap &
\DynamicBabyCandidateOOFOverlap\\
Diginetica & \DynamicDigiCandidatePool & \DynamicDigiDeclaredValidation &
\DynamicDigiTargetsRemoved & \DynamicDigiDeclaredOOFOverlap &
\DynamicDigiCandidateOOFOverlap\\
\bottomrule
\end{tabular}
\end{table}

\begin{table}[h]
\caption{Dataset and protocol statistics.}
\centering\small
\begin{tabular}{lrrrrl}
\toprule
Domain & Items & Train seq. & Validation & Test & Seen-item policy\\
\midrule
Video Games & 25,613 & 94,762 & 5,000 & 94,762 & mask\\
Baby Products & 36,014 & 150,777 & 5,000 & 150,777 & mask\\
Diginetica & 43,098 & 186,670 & 5,000 & 60,858 & retain\\
\bottomrule
\end{tabular}
\end{table}

For cutoff $K$, Recall is
$|\{q:r_q\leq K\}|/|\mathcal Q|$ and nDCG is
$|\mathcal Q|^{-1}\sum_q\mathbf1[r_q\leq K]/\log_2(r_q+1)$.
The allocation objective reported in the paper is
$U=.5\,\mathrm{Recall@6}+.5\,\mathrm{Recall@20}$; Recall@10 and all nDCG
cutoffs are reported in Tables~9--11.

\section{Complete CEARF Memory}
\subsection{Transition evidence}
Only the eight most recent context items are inspected. For context item
$s_a$ at reverse-recency age $a$ and candidate $x$,
\begin{equation}
T_q(x)=\sum_a e^{-.35a}
\frac{\log(1+c(s_a,x))}{\sqrt{1+\mathrm{freq}(x)}}.
\end{equation}
The directed count $c$ sums inverse transition distance within a four-step
window. Each source retains at most 200 outgoing candidates.

\subsection{Neighbour-session evidence}
Let $j$ be an indexed training session. It receives overlap
\begin{equation}
O(q,j)=\sum_a \mathrm{idf}(i_a)e^{-.20a}
\mathbf1[i_a\in j],
\qquad
\mathrm{idf}(i)=\log\!\left(1+
\frac{|\mathcal S|}{|\{j:i\in j\}|}\right).
\end{equation}
The 120 highest-overlap sessions vote for candidates. The vote is normalized
by $\sqrt{|\mathrm{unique}(j)|}$. An item after the most recent matched
position receives direction weight 1; an earlier item receives .25; distance
from the match contributes an additional inverse-square-root factor.

\subsection{Popularity, profiles, and agreement}
Global frequency supplies the third top-120 list and pads a short output.
The declared memory family is:
\begin{center}\small
\begin{tabular}{lrrr}
\toprule
Profile & Transition & Session & Popularity\\
\midrule
transition & 1.00 & 0 & 0\\
session & 0 & 1.00 & 0\\
balanced & .50 & .40 & .10\\
transition--session & .60 & .40 & 0\\
session--transition & .35 & .60 & .05\\
short-safe & .65 & .20 & .15\\
\bottomrule
\end{tabular}
\end{center}
Ranks are combined as $\sum_e w_e/(20+r_e(x))$. If at least two active
memory lists retrieve $x$, its score is multiplied by
$1+.12(v_x-1)$, where $v_x$ is its support count. Short ($L\leq2$) and
longer contexts select profiles independently on the 1,000-query
training-only profile subset. The selected profile is session-only in all
three domains and both regimes; transition and popularity therefore remain
evaluated alternatives/fallbacks rather than hidden contributors to the
primary result.

\section{PASGR Neural Expert}
PASGR expands to \emph{Prototype-Aligned Semantic Graph Retrieval}. It is a
one-layer GRU with 64 hidden units, dropout .2, candidate width 120, and four
training epochs in the dynamic-$\beta$ experiment. For prefix state
$\mathbf h_q$ and item embedding $\mathbf e_i$, exact full-catalogue inference
uses
\[
s_N(q,i)=
\operatorname{norm}(W\,\mathrm{GRU}(q))^\top
\operatorname{norm}(\mathbf e_i).
\]
Video Games item embeddings use the cached 384-dimensional teacher documented
as \texttt{E5-small}. Baby Products uses a
128-dimensional TF--IDF/SVD teacher constructed from available titles,
categories, features, descriptions, and brands; Diginetica uses a
128-dimensional TF--IDF/SVD matrix of product-name tokens. Teacher identity
is resolved from the loaded matrix before evaluation; its shape and SHA-256
prefix are reported directly in the provenance table in Section~10 rather
than inferred from a requested label.

Training contrasts the observed target with 32 negatives drawn from semantic
prototype neighbours, transition-graph neighbours, and random catalogue
items. Cross-entropy on that candidate set is augmented by a recency-weighted
semantic alignment loss and, when enabled, multi-positive in-batch retrieval.
Graph aggregation mixes adjacent-item representations before GRU training;
prototype transport pulls poorly observed items toward semantic centroids.
AdamW uses learning rate $10^{-3}$ and weight decay $10^{-4}$.
The upstream PASGR family crossed graph weight $\{0,.35\}$, prototype
transport on/off, contrastive weight $\{0,.15\}$, and in-batch weight
$\{0,.1\}$. All three frozen cells disable prototype transport, so the
``prototype-aligned'' mechanism is part of the declared PASGR family but not a
hidden contributor to the primary cells.

\begin{table}[h]
\caption{Frozen PASGR cells used by every allocation control. The dynamic
experiment does not reselect them.}
\centering\small
\begin{tabular}{lrrrr}
\toprule
Domain & Graph & Prototype & Contrastive & In-batch\\
\midrule
Video Games & .35 & off & .15 & .10\\
Baby Products & 0 & off & .15 & 0\\
Diginetica & .35 & off & 0 & .10\\
\bottomrule
\end{tabular}
\end{table}

The earlier study evaluated the full 16-cell family on validation. This is an
upstream expert-selection decision, not evidence that dynamic $\beta$ used
validation. A pure
component-attribution study would need to retrain all cells and evaluate
neural-only ranks under the new protocol; we do not relabel the older
fusion-selected grid as such an ablation.

\section{Training-Only Dynamic Allocation}
\subsection{Rank evidence and hard negatives}
For expert $e\in\{M,N\}$ and rank $r_e(x)$,
\begin{equation}
\rr_e(x)=
\begin{cases}
(k+1)/(k+r_e(x)),&x\text{ appears in the top 120},\\
0,&\text{otherwise},
\end{cases}
\qquad k=20.
\end{equation}
Each calibration query uses the 32 non-target union candidates with largest
evidence under either expert, or all candidates when fewer than 32 exist;
item ID breaks ties deterministically. A query is actionable when its target appears
in at least one expert list. Non-actionable queries are retained in coverage
statistics but excluded from optimization because changing $\beta$ cannot
retrieve their missing target from the union.

The serving implementation uses $1/(k+r)$ rather than $(k+1)/(k+r)$.
The omitted common positive factor gives $S_q=F_q/(k+1)$ and therefore exactly
the same ordering. The expert interface is rank-only: each expert enters
allocation through its top-120 ordering, while query features are target-free
at inference. Any strictly monotone raw-score transformation that preserves an
expert's top-120 ordering therefore preserves the OOF training arrays, learned
allocator, and final ranking.

\subsection{Global prior and bounded query residual}
The global model learns one unconstrained logit whose sigmoid is
$\beta_{\mathrm{OOF}}$, initialized at .35. The query model is
\begin{equation}
\beta_q=\beta_{\mathrm{OOF}}+
\Delta_{\mathrm{eff}}\tanh(\mathbf w^\top\widetilde{\mathbf z}_q+b),
\quad
\Delta_{\mathrm{eff}}=\min\{.10,\beta_{\mathrm{OOF}},
1-\beta_{\mathrm{OOF}}\}.
\end{equation}
The primary feature vector contains exactly:
\begin{enumerate}
\item $\log(1+L)$;
\item $\log(1+\mathrm{freq}(i_L))$; and
\item $\mathbf1[i_L\notin\mathcal I_{\mathrm{head}}]$.
\end{enumerate}
Features are computable before recommendation and contain neither target
identity nor a target-hit diagnostic. Mean and standard deviation are fit on
actionable OOF training rows and frozen.

For target $y_q$ and negative $x$, the fused evidence is
$F_q(x)=(1-\beta_q)\rr_M(x)+\beta_q\rr_N(x)$. With temperature $\tau=.10$,
\begin{align}
\mathcal L_{\mathrm{rank}}
&=\frac1{|\mathcal Q|}\sum_q\frac1{|\mathcal H_q|}
\sum_{x\in\mathcal H_q}
\operatorname{softplus}\!\left(
\frac{F_q(x)-F_q(y_q)}{\tau}\right),\\
\mathcal L
&=\mathcal L_{\mathrm{rank}}
+10^{-3}\overline{\beta_q}
+10^{-3}\overline{(\beta_q-\beta_{\mathrm{OOF}})^2}.
\end{align}
The first regularizer is a small neural-admission cost; the second shrinks
query modulation toward the global policy. AdamW uses learning rate .02 for
the 80-epoch scalar fit and $10^{-3}$ for the 50-epoch gate, batch size at
most 512, weight decay $10^{-4}$ for the gate, and deterministic seed order.
The scalar stage minimizes
$\mathcal L_{\mathrm{rank}}+10^{-3}\beta$ with
$\beta_q\equiv\sigma(a)$; only the quadratic residual-prior term is absent.
It is then frozen while the second stage optimizes only $(\mathbf w,b)$ under
the full objective.
For a target--negative pair, define the signed expert-margin advantage
\[
A_{qx}=[\rr_N(y_q)-\rr_N(x)]-[\rr_M(y_q)-\rr_M(x)].
\]
The derivative of the fused target margin with respect to $\beta_q$ is
$A_{qx}$: the gate learns from context when neural evidence has the larger
target-versus-negative margin.

\begin{table}[h]
\caption{Training-only calibration status, mean (SD) over
\DynamicSeedCount{} matched seeds.}
\centering\small
\begin{tabular}{lrrrr}
\toprule
Domain & OOF rows & Actionable & Share & $\beta_{\mathrm{OOF}}$\\
\midrule
Video Games & \DynamicVideoOOFRows & \DynamicVideoActionable &
\DynamicVideoActionableShare & \DynamicVideoGlobalBeta\\
Baby Products & \DynamicBabyOOFRows & \DynamicBabyActionable &
\DynamicBabyActionableShare & \DynamicBabyGlobalBeta\\
Diginetica & \DynamicDigiOOFRows & \DynamicDigiActionable &
\DynamicDigiActionableShare & \DynamicDigiGlobalBeta\\
\bottomrule
\end{tabular}
\end{table}

\subsection{Why the residual is bounded}
The same bounded three-feature linear form, feature set, and hyperparameters
are primary across all domains and seeds; fitted coefficients vary by
replicate. The direct
16-feature MLP is a capacity stress test. Its gains are small and
seed-dependent; on Diginetica, three of four MLP variants are lower than the
primary and the remaining no-memory-certainty MLP is higher by only .00025
$U$. The primary gate
treats $\beta_{\mathrm{OOF}}$ as a learned prior and permits only a $\pm.10$
correction. This does not guarantee improvement or bound ranking loss; it
limits only per-item fused-score perturbation. The linear inside the $\tanh$
has four optimized coefficients. Together with the scalar prior, the allocator
has five optimized coefficients; six feature means/standard deviations are
frozen preprocessing state, not optimized parameters.

\paragraph{Bounded expert-disagreement certificate.}
Analytically, finite real logits give $\beta_q\in(0,1)$ because
$\Delta_{\mathrm{eff}}\leq
\min(\beta_{\mathrm{OOF}},1-\beta_{\mathrm{OOF}})$ and
$|\tanh(\cdot)|<1$; the implementation-safe guarantee is $[0,1]$, and all
persisted allocations are strictly interior. Moreover, with
$F_{\mathrm{OOF}}(x)=(1-\beta_{\mathrm{OOF}})\rr_M(x)+
\beta_{\mathrm{OOF}}\rr_N(x)$,
\[
|F_q(x)-F_{\mathrm{OOF}}(x)|
=|\beta_q-\beta_{\mathrm{OOF}}|\,|\rr_N(x)-\rr_M(x)|
\leq\Delta_{\mathrm{eff}}\leq.10,
\]
since each reciprocal-rank evidence lies in $[0,1]$. More sharply, let
$\delta_q=\beta_q-\beta_{\mathrm{OOF}}$,
$D_q(x)=\rr_N(x)-\rr_M(x)$, and
$m_{\mathrm{OOF}}(a,b)=F_{\mathrm{OOF}}(a)-F_{\mathrm{OOF}}(b)$. Then
\[
F_q(a)-F_q(b)=m_{\mathrm{OOF}}(a,b)+
\delta_q\{D_q(a)-D_q(b)\}.
\]
The baseline order is certified whenever
$|m_{\mathrm{OOF}}(a,b)|>
|\delta_q|\,|D_q(a)-D_q(b)|$, and conservatively whenever
$|m_{\mathrm{OOF}}(a,b)|>2\Delta_{\mathrm{eff}}$. If
$D_q(a)=D_q(b)$, the pair is allocation-invariant. For a target--negative
pair, $D_q(y_q)-D_q(x)=A_{qx}$ and
\[
\frac{\partial\ell_{qx}}{\partial\beta_q}
=-\frac{A_{qx}}{\tau}
\sigma\!\left(\frac{F_q(x)-F_q(y_q)}{\tau}\right).
\]
Thus the same expert-disagreement term determines the learning signal and the
maximum permitted pair intervention. This is a score-margin certificate, not
a Recall/nDCG guarantee.
Conditional on the
declared source split and frozen
experts, $\beta_{\mathrm{OOF}},\mathbf w,b$ are functions only of OOF
training ranks, OOF training targets, and training-derived feature
statistics. Declared-validation and test targets therefore cannot affect the
allocator parameters. This construction establishes allocation provenance;
it does not by itself guarantee test improvement.

\section{Out-of-Fit Split and Leakage Controls}
Eligible source sessions have at least three events and do not overlap the
declared 5,000-query validation subset. Their identifiers are ordered by a
stable BLAKE2 fraction. At most 5,000 (10\% when smaller) are held out:
\begin{itemize}
\item the first 1,000 become profile-lock leave-last-out queries;
\item the remaining 4,000 become gate-calibration leave-last-out queries;
\item all 5,000 source sessions are entirely absent from the inner CEARF
index and PASGR training fit.
\end{itemize}
Consequently, the target event, its incoming transition, and every earlier
event from the same source are absent from the experts that predict it. The
profile/gate fingerprints are SHA-256 hashes computed before
declared-validation/test prediction.

On Diginetica, declared-validation target events and incoming transitions are
removed from tuning sessions; on Amazon those targets are already outside the
training histories. The complete training split is restored only for the
single final test expert fit. Declared-validation labels are allowed only for
reporting the frozen system, never as arguments to either beta fitter. We
separately record candidate-pool size, declared-subset size, and zero source
overlap between the declared subset and OOF calibration.

\begin{table}[h]
\caption{Fixed split identities. SHA-256 fingerprints are shown as stable
abbreviated identity checks.}
\centering\small
\begin{tabular}{lrrll}
\toprule
Domain & Eligible & Held & Profile hash & Gate hash\\
\midrule
Video Games & \DynamicVideoEligible & 5,000 &
\DynamicVideoProfileHash\ldots & \DynamicVideoGateHash\ldots\\
Baby Products & \DynamicBabyEligible & 5,000 &
\DynamicBabyProfileHash\ldots & \DynamicBabyGateHash\ldots\\
Diginetica & \DynamicDigiEligible & 5,000 &
\DynamicDigiProfileHash\ldots & \DynamicDigiGateHash\ldots\\
\bottomrule
\end{tabular}
\end{table}

\section{Complete Allocation Results}
% Generated by summarize_dynamic_beta.py; do not edit manually.
\begin{table}[H]
\caption{Allocation ablation on full-catalogue test queries. Values are mean $\pm$ SD over matched seeds; $U=.5R@6+.5R@20$. No row uses validation to fit or select $\beta$.}
\label{tab:allocation}
\centering\scriptsize
\setlength{\tabcolsep}{2.7pt}
\resizebox{\textwidth}{!}{%
\begin{tabular}{llrrrr}
\toprule
Domain & Allocation & R@6 & R@20 & nD@20 & $U$\\
\midrule
Video & Memory endpoint & .06932 (.00000) & .11901 (.00000) & .06009 (.00000) & .09417 (.00000)\\
 & Neural endpoint & .06158 (.00019) & .12571 (.00007) & .05611 (.00007) & .09364 (.00012)\\
 & Fixed $.5$ & .07961 (.00022) & .14610 (.00018) & .06937 (.00025) & .11285 (.00009)\\
 & OOF global & .07799 (.00036) & .14719 (.00033) & .06886 (.00011) & .11259 (.00027)\\
 & OOF short/long & .07830 (.00049) & .14726 (.00014) & .06900 (.00015) & .11278 (.00023)\\
 & \textbf{Dynamic $\beta_q$} & .07827 (.00035) & .14735 (.00020) & .06904 (.00015) & .11281 (.00021)\\
\midrule
Baby & Memory endpoint & .02288 (.00000) & .03879 (.00000) & .02017 (.00000) & .03084 (.00000)\\
 & Neural endpoint & .02337 (.00060) & .04873 (.00050) & .02164 (.00040) & .03605 (.00049)\\
 & Fixed $.5$ & .03041 (.00038) & .05555 (.00038) & .02666 (.00017) & .04298 (.00031)\\
 & OOF global & .03011 (.00031) & .05671 (.00042) & .02690 (.00015) & .04341 (.00031)\\
 & OOF short/long & .03038 (.00047) & .05655 (.00053) & .02690 (.00019) & .04346 (.00040)\\
 & \textbf{Dynamic $\beta_q$} & .03017 (.00037) & .05669 (.00037) & .02691 (.00016) & .04343 (.00032)\\
\midrule
Diginetica & Memory endpoint & .28941 (.00000) & .49428 (.00000) & .24417 (.00000) & .39185 (.00000)\\
 & Neural endpoint & .26399 (.00212) & .44852 (.00302) & .22840 (.00124) & .35626 (.00249)\\
 & Fixed $.5$ & .30866 (.00050) & .51430 (.00135) & .25658 (.00070) & .41148 (.00087)\\
 & OOF global & .30958 (.00105) & .51702 (.00142) & .25786 (.00056) & .41330 (.00103)\\
 & OOF short/long & .30926 (.00115) & .51692 (.00159) & .25775 (.00059) & .41309 (.00117)\\
 & \textbf{Dynamic $\beta_q$} & .30964 (.00135) & .51701 (.00142) & .25792 (.00052) & .41332 (.00115)\\
\bottomrule
\end{tabular}
}
\end{table}

\begin{table}[H]
\caption{Dynamic $\beta_q$ minus training-only global $\beta$; paired query-level 95\% bootstrap intervals after averaging matched-seed outcomes.}
\label{tab:dynamic-delta}
\centering\scriptsize
\setlength{\tabcolsep}{3.5pt}
\resizebox{\textwidth}{!}{%
\begin{tabular}{lrrrr}
\toprule
Domain & $\Delta$R@6 & $\Delta$R@20 & $\Delta$nD@20 & $\Delta U$\\
\midrule
Video & +.00028 [+.00011,+.00044] & +.00016 [-.00003,+.00036] & +.00018 [+.00012,+.00024] & +.00022 [+.00009,+.00035]\\
Baby & +.00005 [-.00001,+.00011] & -.00002 [-.00009,+.00004] & +.00000 [-.00002,+.00002] & +.00002 [-.00003,+.00006]\\
Diginetica & +.00005 [-.00018,+.00030] & -.00002 [-.00024,+.00020] & +.00006 [-.00001,+.00013] & +.00002 [-.00015,+.00018]\\
\bottomrule
\end{tabular}
}
\end{table}

% Generated by summarize_dynamic_beta.py; do not edit manually.
\begin{table}[H]
\caption{Complete allocation metrics on full-catalogue test queries. Values are mean (SD) over matched seeds.}
\label{tab:full-allocation-metrics}
\centering\scriptsize
\resizebox{\textwidth}{!}{%
\begin{tabular}{llrrrrrrr}
\toprule
Domain & Allocation & R@6 & nD@6 & R@10 & nD@10 & R@20 & nD@20 & $U$\\
\midrule
Video & Fixed $.5$ & .07961 (.00022) & .05110 (.00022) & .10505 (.00034) & .05901 (.00019) & .14610 (.00018) & .06937 (.00025) & .11285 (.00009)\\
 & OOF global & .07799 (.00036) & .04992 (.00020) & .10347 (.00024) & .05785 (.00024) & .14719 (.00033) & .06886 (.00011) & .11259 (.00027)\\
 & OOF short/long & .07830 (.00049) & .05012 (.00027) & .10403 (.00031) & .05811 (.00025) & .14726 (.00014) & .06900 (.00015) & .11278 (.00023)\\
 & Dynamic $\beta_q$ & .07827 (.00035) & .05012 (.00024) & .10388 (.00018) & .05809 (.00024) & .14735 (.00020) & .06904 (.00015) & .11281 (.00021)\\
\midrule
Baby & Fixed $.5$ & .03041 (.00038) & .01978 (.00023) & .03959 (.00028) & .02265 (.00018) & .05555 (.00038) & .02666 (.00017) & .04298 (.00031)\\
 & OOF global & .03011 (.00031) & .01963 (.00020) & .03992 (.00029) & .02268 (.00020) & .05671 (.00042) & .02690 (.00015) & .04341 (.00031)\\
 & OOF short/long & .03038 (.00047) & .01976 (.00025) & .03995 (.00022) & .02273 (.00017) & .05655 (.00053) & .02690 (.00019) & .04346 (.00040)\\
 & Dynamic $\beta_q$ & .03017 (.00037) & .01966 (.00022) & .03997 (.00031) & .02270 (.00020) & .05669 (.00037) & .02691 (.00016) & .04343 (.00032)\\
\midrule
Diginetica & Fixed $.5$ & .30866 (.00050) & .19970 (.00049) & .39255 (.00114) & .22583 (.00068) & .51430 (.00135) & .25658 (.00070) & .41148 (.00087)\\
 & OOF global & .30958 (.00105) & .20043 (.00043) & .39470 (.00078) & .22696 (.00038) & .51702 (.00142) & .25786 (.00056) & .41330 (.00103)\\
 & OOF short/long & .30926 (.00115) & .20028 (.00047) & .39427 (.00084) & .22677 (.00041) & .51692 (.00159) & .25775 (.00059) & .41309 (.00117)\\
 & Dynamic $\beta_q$ & .30964 (.00135) & .20051 (.00057) & .39463 (.00081) & .22700 (.00039) & .51701 (.00142) & .25792 (.00052) & .41332 (.00115)\\
\bottomrule
\end{tabular}
}
\end{table}

The first table includes the two endpoints and every principal allocation
control; the second gives all declared cutoffs for the non-endpoint allocation
policies. Complete per-cutoff results are reported in Tables~9--11.

\subsection{Per-query rescues and dynamic range}
\begin{table}[h]
\caption{Dynamic behavior on full test sets, mean (SD) over matched seeds.
A rescue is an R@20 miss by memory changed to a hit; damage is the reverse.}
\centering\small
\begin{tabular}{lrrrrr}
\toprule
Domain & Mean $\beta_q$ & SD & Rescues & Damage & Net\\
\midrule
Video Games & \DynamicVideoBetaMean & \DynamicVideoBetaWithinSD &
\DynamicVideoRescues & \DynamicVideoDamage & \DynamicVideoNetRescues\\
Baby Products & \DynamicBabyBetaMean & \DynamicBabyBetaWithinSD &
\DynamicBabyRescues & \DynamicBabyDamage & \DynamicBabyNetRescues\\
Diginetica & \DynamicDigiBetaMean & \DynamicDigiBetaWithinSD &
\DynamicDigiRescues & \DynamicDigiDamage & \DynamicDigiNetRescues\\
\bottomrule
\end{tabular}
\end{table}
The large net counts are mainly a fusion effect. To isolate dynamic
allocation, compare with the global policy: dynamic changes mean net rescues
by \DynamicVideoNetVsGlobal{} on Video,
\DynamicBabyNetVsGlobal{} on Baby, and
\DynamicDigiNetVsGlobal{} on Diginetica.

% Generated by audit_short_paper_baselines.py; do not hand-edit.
\section{Audited External-Baseline Details}
\noindent The following tables report the full-catalogue baseline results used in the main-paper comparison. A dagger marks deterministic methods: their rows for seeds 42/123/456 are identical copies for schema compatibility, not independent stochastic fits.  STAN uses loader session order as its recency coordinate because original timestamps are unavailable.  Baseline definitions follow V-SKNN \cite{ludewig2018vsknn}, STAN \cite{garg2019stan}, GRU4Rec \cite{hidasi2016gru4rec}, NARM \cite{li2017narm}, SR-GNN \cite{wu2019srgnn}, and SIGMA \cite{sigma2025}; the last is our compatible reimplementation.

\begin{table}[h]
\caption{Validation grids and selected neighborhood configurations.}
\centering\scriptsize
\begin{tabular}{llp{.58\textwidth}}
\toprule
Domain & Method & Grid and selected configuration\\
\midrule
Video Games & V-SKNN & 8 cells: $k\in\{100,500\}$, sample $\in\{1000,5000\}$, weighting $\in\{\mathrm{div},\mathrm{quadratic}\}$; selected $k=500$, sample $=5000$, div, exclude-seen=true.\\
Video Games & STAN & 16 cells: $k\in\{100,500\}$, sample $\in\{1000,5000\}$, $\lambda_{spw}\in\{1.02,2.0\}$, $\lambda_{snh}\in\{\emptyset,5000\}$; selected $k=500$, sample $=5000$, $\lambda_{spw}=2.0$, $\lambda_{snh}=\emptyset$, exclude-seen=true.\\
Baby Products & V-SKNN & 8 cells: $k\in\{100,500\}$, sample $\in\{1000,5000\}$, weighting $\in\{\mathrm{div},\mathrm{quadratic}\}$; selected $k=500$, sample $=5000$, div, exclude-seen=true.\\
Baby Products & STAN & 16 cells: $k\in\{100,500\}$, sample $\in\{1000,5000\}$, $\lambda_{spw}\in\{1.02,2.0\}$, $\lambda_{snh}\in\{\emptyset,5000\}$; selected $k=500$, sample $=5000$, $\lambda_{spw}=2.0$, $\lambda_{snh}=\emptyset$, exclude-seen=true.\\
Diginetica & V-SKNN & 8 cells: $k\in\{100,500\}$, sample $\in\{1000,5000\}$, weighting $\in\{\mathrm{div},\mathrm{quadratic}\}$; selected $k=100$, sample $=5000$, quadratic, exclude-seen=false.\\
Diginetica & STAN & 16 cells: $k\in\{100,500\}$, sample $\in\{1000,5000\}$, $\lambda_{spw}\in\{1.02,2.0\}$, $\lambda_{snh}\in\{\emptyset,5000\}$; selected $k=500$, sample $=5000$, $\lambda_{spw}=1.02$, $\lambda_{snh}=\emptyset$, exclude-seen=false.\\
\bottomrule
\end{tabular}
\par\smallskip\scriptsize Selection maximizes $0.5R@6+0.5R@20$ on 5,000 declared validation queries. Score weighting is \texttt{div}; $\lambda_{inh}=2.05$ is fixed.
\end{table}

\begin{table}[h]
\caption{Compact full-catalog baseline metrics on Video Games. Neural entries are mean $\pm$ sample SD over seeds 42/123/456. R@10 and earlier-cutoff nDCG are omitted for compactness.}
\centering\scriptsize
\begin{tabular}{lrrr}
\toprule
Method & R@6 & R@20 & nDCG@20\\
\midrule
Transition-only$^\dagger$ & 0.01712 & 0.04223 & 0.01733\\
V-SKNN$^\dagger$ & 0.06713 & 0.11936 & 0.05852\\
STAN$^\dagger$ & 0.06841 & 0.12115 & 0.05915\\
GRU4Rec & 0.05238 $\pm$ 0.00113 & 0.10368 $\pm$ 0.00089 & 0.04733 $\pm$ 0.00088\\
NARM & 0.06956 $\pm$ 0.00069 & 0.13705 $\pm$ 0.00081 & 0.06316 $\pm$ 0.00037\\
SR-GNN & 0.06154 $\pm$ 0.00037 & 0.12267 $\pm$ 0.00060 & 0.05558 $\pm$ 0.00031\\
SIGMA-compatible & 0.04367 $\pm$ 0.00104 & 0.08644 $\pm$ 0.00215 & 0.03970 $\pm$ 0.00077\\
\bottomrule
\end{tabular}
\end{table}

\begin{table}[h]
\caption{Compact full-catalog baseline metrics on Baby Products. Neural entries are mean $\pm$ sample SD over seeds 42/123/456. R@10 and earlier-cutoff nDCG are omitted for compactness.}
\centering\scriptsize
\begin{tabular}{lrrr}
\toprule
Method & R@6 & R@20 & nDCG@20\\
\midrule
Transition-only$^\dagger$ & 0.00299 & 0.00735 & 0.00311\\
V-SKNN$^\dagger$ & 0.02604 & 0.05038 & 0.02388\\
STAN$^\dagger$ & 0.02721 & 0.04944 & 0.02425\\
GRU4Rec & 0.01939 $\pm$ 0.00038 & 0.04501 $\pm$ 0.00056 & 0.01882 $\pm$ 0.00019\\
NARM & 0.01195 $\pm$ 0.00060 & 0.02990 $\pm$ 0.00018 & 0.01117 $\pm$ 0.00006\\
SR-GNN & 0.02332 $\pm$ 0.00038 & 0.05208 $\pm$ 0.00127 & 0.02230 $\pm$ 0.00032\\
SIGMA-compatible & 0.00856 $\pm$ 0.00096 & 0.02473 $\pm$ 0.00059 & 0.00897 $\pm$ 0.00039\\
\bottomrule
\end{tabular}
\end{table}

\begin{table}[h]
\caption{Compact full-catalog baseline metrics on Diginetica. Neural entries are mean $\pm$ sample SD over seeds 42/123/456. R@10 and earlier-cutoff nDCG are omitted for compactness.}
\centering\scriptsize
\begin{tabular}{lrrr}
\toprule
Method & R@6 & R@20 & nDCG@20\\
\midrule
Transition-only$^\dagger$ & 0.18392 & 0.40102 & 0.16687\\
V-SKNN$^\dagger$ & 0.30481 & 0.50506 & 0.25184\\
STAN$^\dagger$ & 0.31785 & 0.51502 & 0.25952\\
GRU4Rec & 0.23008 $\pm$ 0.00065 & 0.41785 $\pm$ 0.00108 & 0.19557 $\pm$ 0.00030\\
NARM & 0.31341 $\pm$ 0.00135 & 0.53406 $\pm$ 0.00052 & 0.26321 $\pm$ 0.00037\\
SR-GNN & 0.27638 $\pm$ 0.00256 & 0.49267 $\pm$ 0.00468 & 0.23599 $\pm$ 0.00237\\
SIGMA-compatible & 0.20528 $\pm$ 0.00589 & 0.37219 $\pm$ 0.01094 & 0.17374 $\pm$ 0.00489\\
\bottomrule
\end{tabular}
\end{table}

\begin{table}[h]
\caption{Validation-selected epochs for ID-only neural baselines.}
\centering\scriptsize
\begin{tabular}{lrrrr}
\toprule
Domain & GRU4Rec & NARM & SR-GNN & SIGMA-compatible\\
\midrule
Video Games & 12/12/12 & 12/12/12 & 5/5/5 & 6/6/6\\
Baby Products & 9/8/10 & 1/1/3 & 5/5/4 & 1/2/1\\
Diginetica & 12/12/12 & 11/12/11 & 4/4/5 & 7/9/10\\
\bottomrule
\end{tabular}
\par\smallskip\scriptsize Shared fixed settings: 64 dimensions, 50-item maximum context, full-catalog cross-entropy, AdamW, weight decay $10^{-5}$, gradient clipping at 5, batch 512 (SR-GNN: 128), learning rate $10^{-3}$ (SIGMA-compatible: $5\!\times\!10^{-4}$).  Only epoch is validation-selected; these are our local reimplementations, and SIGMA-compatible is not official SIGMA.
\end{table}

\begin{table}[h]
\caption{Resolved semantic-teacher provenance. Reported SHA-256 prefixes identify the cached matrices used.}
\centering\small
\begin{tabular}{lp{.40\textwidth}ll}
\toprule
Domain & Resolved teacher & Shape & SHA-256\\
\midrule
Video Games & E5-small cache & $25613\times384$ & \texttt{a22a31407461}\\
Baby Products & TF-IDF/SVD cache & $36014\times128$ & \texttt{3191943cc015}\\
Diginetica & TF-IDF/SVD product-name cache & $43098\times128$ & \texttt{e31c0d501b42}\\
\bottomrule
\end{tabular}
\end{table}

\begin{table}[h]
\caption{Matched-teacher NARM final-cutoff metrics on Amazon domains. Values are mean $\pm$ sample SD over seeds 42/123/456; no matched-teacher Diginetica run is retained.}
\centering\small
\begin{tabular}{lrr}
\toprule
Domain & R@20 & nDCG@20\\
\midrule
Video Games & 0.15387 $\pm$ 0.00051 & 0.07112 $\pm$ 0.00029\\
Baby Products & 0.06628 $\pm$ 0.00014 & 0.02914 $\pm$ 0.00010\\
\bottomrule
\end{tabular}
\par\smallskip\scriptsize Only the final cutoff is shown because the matched-teacher claim concerns R@20 and nDCG@20.
\par\smallskip\scriptsize The earlier fairness record mislabels the Video Games teacher: resolving the loaded matrix gives the higher-priority E5-small cache before TF--IDF/SVD. The Video cache has no encoder sidecar; the reported SHA-256 prefix identifies the matrix used. Exact matched-teacher NARM R@20 means are .153869695 and .066278898.
\end{table}

\noindent\textbf{No-metadata scope.} The available no-metadata results do not ablate the current training-only dynamic-$\beta$ method; they belong to the earlier validation-selected short/long-$\beta$ family. Only the nested Diginetica result is protocol-safe as a static diagnostic.


\FloatBarrier

\subsection{Neural-expert sensitivity}
This completed 3-domain, 3-seed control replaces PASGR by ID-only NARM while
holding CEARF memory, OOF split, continuous objective, bounded gate, RRF
$k=20$, and full-catalogue test protocol fixed. It is not metadata-matched on
Amazon because PASGR uses semantic initialization; its role is to show that the
allocator can absorb a stronger ID-only neural expert. The swap closes the
Diginetica external-baseline gap (dynamic R@20 .53939 versus NARM .53406) but
reduces Baby relative to the PASGR primary, consistent with Baby's metadata
benefit.
% Generated by summarize_dynamic_beta_expert_swap.py.
\begin{table}[H]
\caption{ID-only NARM expert swap. Values are test R@20 mean (SD) over three completed seeds. The CEARF memory, OOF split, continuous objective, bounded gate, and $k=20$ are held fixed; the test expert is freshly refit on complete training sessions.}
\label{tab:narm-swap}
\centering\small
\setlength{\tabcolsep}{4pt}
\begin{tabular}{llrrr}
\toprule
Domain & Neural expert & Neural-only & OOF global & Dynamic $\beta_q$\\
\midrule
Video & NARM (ID-only) & .13763 (.00088) & .14819 (.00041) & .14854 (.00034)\\
\midrule
Baby & NARM (ID-only) & .03880 (.00785) & .04994 (.00542) & .05001 (.00534)\\
\midrule
Diginetica & NARM (ID-only) & .53178 (.00025) & .53910 (.00084) & .53939 (.00061)\\
\bottomrule
\end{tabular}
\end{table}

\FloatBarrier

\section{Gate-Capacity Ablation}
The implementation computes 16 target-free diagnostics for ablation:
\begin{table}[H]
\caption{Target-free diagnostics available to gate-capacity ablations. The
primary gate uses only features 1--3.}
\label{tab:gate-diagnostics}
\centering\small
\begin{tabular}{p{.04\textwidth}p{.39\textwidth}p{.04\textwidth}p{.39\textwidth}}
\toprule
ID & Feature & ID & Feature\\
\midrule
1 & log context length & 9 & memory top-1 in neural top-5\\
2 & log last-item training frequency & 10 & neural top-1 in memory top-5\\
3 & last-item tail indicator & 11 & memory top-1 component support\\
4 & maximum memory agreement at 5 & 12 & neural top-1 component support\\
5 & memory/neural top-1 equality & 13 & mean memory Jaccard at 10\\
6 & memory--neural Jaccard at 5 & 14 & neural novelty outside memory top-120\\
7 & memory--neural Jaccard at 20 & 15 & log transition branching factor\\
8 & reciprocal-rank overlap at 20 & 16 & dominant transition share\\
\bottomrule
\end{tabular}
\end{table}

\clearpage
\begin{table}[h]
\caption{Gate-capacity ablation, test $U$ mean (SD) over matched seeds.
``Primary'' is features 1--3, linear, bounded by .10.}
\centering\small
\setlength{\tabcolsep}{4pt}
\resizebox{\textwidth}{!}{%
\begin{tabular}{lrrr}
\toprule
Gate & Video & Baby & Diginetica\\
\midrule
Primary context linear & \DynamicVideoPrimaryU & \DynamicBabyPrimaryU &
\DynamicDigiPrimaryU\\
Context MLP (16 hidden) & \DynamicVideoContextMLPU &
\DynamicBabyContextMLPU & \DynamicDigiContextMLPU\\
All-feature linear & \DynamicVideoFullLinearU & \DynamicBabyFullLinearU &
\DynamicDigiFullLinearU\\
All-feature MLP & \DynamicVideoFullMLPU & \DynamicBabyFullMLPU &
\DynamicDigiFullMLPU\\
MLP without cross-expert features & \DynamicVideoNoCrossExpertU &
\DynamicBabyNoCrossExpertU & \DynamicDigiNoCrossExpertU\\
MLP without memory-certainty features &
\DynamicVideoNoMemoryCertaintyU & \DynamicBabyNoMemoryCertaintyU &
\DynamicDigiNoMemoryCertaintyU\\
\bottomrule
\end{tabular}
}
\end{table}
No domain-specific post-hoc architecture is selected. The all-feature linear
gate has a slightly higher mean in each domain but loses to the primary on one
matched seed in each. On Diginetica, three of four MLP variants are lower than
the primary and the remaining no-memory-certainty MLP is higher by only .00025
$U$. The same
three-feature linear form and hyperparameters therefore remain primary across
domains and seeds.

\section{Fusion Sensitivity}
\begin{table}[h]
\caption{Dynamic test $U$ mean (SD) after refitting the continuous global
prior and gate on OOF training ranks for each RRF constant.}
\centering\small
\begin{tabular}{lrrr}
\toprule
Constant & Video & Baby & Diginetica\\
\midrule
$k=10$ & \DynamicVideoKTenU & \DynamicBabyKTenU & \DynamicDigiKTenU\\
$k=20$ & \DynamicVideoKTwentyU & \DynamicBabyKTwentyU &
\DynamicDigiKTwentyU\\
$k=60$ & \DynamicVideoKSixtyU & \DynamicBabyKSixtyU &
\DynamicDigiKSixtyU\\
\bottomrule
\end{tabular}
\end{table}
This is a sensitivity test, not a validation search: primary $k=20$ is held
fixed across domains and seeds rather than reselected from this table. Raw
reciprocal ranks are never relabelled as raw
expert scores. A normalized-score operator control is reported only if frozen
checkpoint inference exactly reconstructs the persisted top-120 expert IDs.
CombSUM min--max normalizes each expert per query over the union of those
persisted top-120 candidates (absent candidates receive zero). It is evaluated
under frozen $\beta_q$ and parameter-free $\beta=.5$; neither allocation is
reoptimized for CombSUM. This is not a full-catalogue score-normalization
claim.

% Generated by summarize_dynamic_beta_fusion_control.py.
\begin{table}[h]
\caption{Fusion-operator perturbation under the frozen dynamic $\beta_q$ and parameter-free $\beta=.5$. Values are mean (SD) over matched seeds. Intervals are paired query-level 95\% bootstrap intervals for normalized CombSUM minus weighted RRF.}
\label{tab:fusion-operator}
\centering\scriptsize
\setlength{\tabcolsep}{2pt}
\begin{tabular}{llrrr}
\toprule
Domain & Allocation & RRF R@20 & CombSUM R@20 & $\Delta$R@20 [95\% CI]\\
\midrule
Video & Dynamic $\beta_q$ & .14735 (.00020) & .14602 (.00058) & -.00133 [-.00175,-.00090]\\
Video & $\beta=.5$ & .14610 (.00018) & .14378 (.00042) & -.00232 [-.00275,-.00190]\\
Baby & Dynamic $\beta_q$ & .05669 (.00037) & .05534 (.00052) & -.00134 [-.00159,-.00110]\\
Baby & $\beta=.5$ & .05555 (.00038) & .05410 (.00038) & -.00146 [-.00169,-.00122]\\
Diginetica & Dynamic $\beta_q$ & .51701 (.00142) & .51429 (.00018) & -.00272 [-.00367,-.00176]\\
Diginetica & $\beta=.5$ & .51430 (.00135) & .51625 (.00068) & +.00194 [+.00095,+.00296]\\
\bottomrule
\end{tabular}
\par\smallskip\scriptsize The table reports the R@20 control used for the fusion-operator claim; no claim depends on omitted cutoffs.
\end{table}


With the frozen dynamic allocation, weighted RRF improves R@20 over normalized
CombSUM by .00133, .00134, and .00272 on Video, Baby, and Diginetica,
respectively. The equal-allocation Diginetica result reverses in favor of
CombSUM. Hence the control supports rank-space RRF for the proposed allocator,
but it also shows an operator--allocation interaction; it does not establish
unconditional RRF superiority.

\section{Allocation Controls}
% Auto-generated by run_dynamic_beta_allocation_controls.py.
% Values are means over completed seeds unless stated otherwise.
\begingroup\small
\setlength{\tabcolsep}{3pt}
\begin{longtable}{llrrrr}
\caption{Training-only allocation controls across all domains. Learned policies use frozen OOF training ranks; reassignment is a fixed target-free control, and no validation or test labels enter fitting or selection.}\label{tab:allocation-controls-all}\\
\toprule
Domain & Policy & R@6 & R@20 & nDCG@20 & $U$ \\
\midrule
\endfirsthead
\multicolumn{6}{c}{\tablename\ \thetable{} -- continued}\\
\toprule
Domain & Policy & R@6 & R@20 & nDCG@20 & $U$ \\
\midrule
\endhead
Video Games & OOF global & 0.07799 & 0.14719 & 0.06886 & 0.11259 \\
Video Games & Head/tail scalars & 0.07821 & 0.14730 & 0.06896 & 0.11275 \\
Video Games & Short/long $\times$ head/tail & 0.07827 & 0.14728 & 0.06901 & 0.11278 \\
Video Games & Dynamic $\Delta=.05$ & 0.07820 & 0.14729 & 0.06900 & 0.11274 \\
Video Games & Dynamic $\Delta=.10$ (primary) & 0.07827 & 0.14735 & 0.06904 & 0.11281 \\
Video Games & Primary $\beta_q$ reassigned & 0.07817 & 0.14725 & 0.06890 & 0.11271 \\
Video Games & Dynamic $\Delta=.20$ & 0.07840 & 0.14744 & 0.06912 & 0.11292 \\
Video Games & Length only & 0.07817 & 0.14720 & 0.06894 & 0.11269 \\
Video Games & Frequency only & 0.07819 & 0.14728 & 0.06896 & 0.11274 \\
Video Games & Tail only & 0.07816 & 0.14731 & 0.06895 & 0.11274 \\
Video Games & Drop length & 0.07819 & 0.14733 & 0.06898 & 0.11276 \\
Video Games & Drop frequency & 0.07824 & 0.14728 & 0.06901 & 0.11276 \\
Video Games & Drop tail & 0.07825 & 0.14734 & 0.06904 & 0.11279 \\
Video Games & No admission cost & 0.07828 & 0.14734 & 0.06904 & 0.11281 \\
Video Games & No prior penalty & 0.07827 & 0.14735 & 0.06904 & 0.11281 \\
\midrule
Baby Products & OOF global & 0.03011 & 0.05671 & 0.02690 & 0.04341 \\
Baby Products & Head/tail scalars & 0.03018 & 0.05661 & 0.02688 & 0.04339 \\
Baby Products & Short/long $\times$ head/tail & 0.03018 & 0.05643 & 0.02684 & 0.04331 \\
Baby Products & Dynamic $\Delta=.05$ & 0.03014 & 0.05672 & 0.02691 & 0.04343 \\
Baby Products & Dynamic $\Delta=.10$ (primary) & 0.03017 & 0.05669 & 0.02691 & 0.04343 \\
Baby Products & Primary $\beta_q$ reassigned & 0.03018 & 0.05671 & 0.02691 & 0.04345 \\
Baby Products & Dynamic $\Delta=.20$ & 0.03016 & 0.05670 & 0.02690 & 0.04343 \\
Baby Products & Length only & 0.03015 & 0.05670 & 0.02690 & 0.04343 \\
Baby Products & Frequency only & 0.03015 & 0.05666 & 0.02690 & 0.04340 \\
Baby Products & Tail only & 0.03017 & 0.05664 & 0.02689 & 0.04340 \\
Baby Products & Drop length & 0.03016 & 0.05667 & 0.02690 & 0.04341 \\
Baby Products & Drop frequency & 0.03017 & 0.05668 & 0.02690 & 0.04342 \\
Baby Products & Drop tail & 0.03016 & 0.05671 & 0.02691 & 0.04343 \\
Baby Products & No admission cost & 0.03017 & 0.05669 & 0.02691 & 0.04343 \\
Baby Products & No prior penalty & 0.03017 & 0.05669 & 0.02691 & 0.04343 \\
\midrule
Diginetica & OOF global & 0.30958 & 0.51702 & 0.25786 & 0.41330 \\
Diginetica & Head/tail scalars & 0.30984 & 0.51707 & 0.25798 & 0.41346 \\
Diginetica & Short/long $\times$ head/tail & 0.30965 & 0.51695 & 0.25792 & 0.41330 \\
Diginetica & Dynamic $\Delta=.05$ & 0.30961 & 0.51704 & 0.25791 & 0.41333 \\
Diginetica & Dynamic $\Delta=.10$ (primary) & 0.30964 & 0.51701 & 0.25792 & 0.41332 \\
Diginetica & Primary $\beta_q$ reassigned & 0.30960 & 0.51685 & 0.25784 & 0.41322 \\
Diginetica & Dynamic $\Delta=.20$ & 0.30956 & 0.51702 & 0.25791 & 0.41329 \\
Diginetica & Length only & 0.30947 & 0.51692 & 0.25781 & 0.41320 \\
Diginetica & Frequency only & 0.30961 & 0.51706 & 0.25791 & 0.41334 \\
Diginetica & Tail only & 0.30985 & 0.51709 & 0.25794 & 0.41347 \\
Diginetica & Drop length & 0.30974 & 0.51710 & 0.25795 & 0.41342 \\
Diginetica & Drop frequency & 0.30966 & 0.51697 & 0.25789 & 0.41331 \\
Diginetica & Drop tail & 0.30956 & 0.51698 & 0.25786 & 0.41327 \\
Diginetica & No admission cost & 0.30959 & 0.51701 & 0.25791 & 0.41330 \\
Diginetica & No prior penalty & 0.30964 & 0.51701 & 0.25792 & 0.41332 \\
\bottomrule
\end{longtable}
\noindent R@10, nDCG@6, and nDCG@10 are omitted because no allocation-control claim depends on them.
\endgroup

\begin{table}[H]
\centering
\small
\caption{Learned assignment versus fixed reassignment across all domains. Entries are primary minus reassigned after matched-seed averaging; paired query-level 95\% bootstrap intervals are unadjusted across metrics.}
\label{tab:beta-assignment-control}
\begin{tabular}{llrr}
\toprule
Domain & Metric & Difference & 95\% CI\\
\midrule
Video Games & R@20 & +0.00010 & [-0.00011, +0.00031] \\
Video Games & nDCG@6 & +0.00013 & [+0.00006, +0.00021] \\
Video Games & nDCG@10 & +0.00015 & [+0.00007, +0.00022] \\
Video Games & nDCG@20 & +0.00015 & [+0.00008, +0.00021] \\
Video Games & $U$ & +0.00010 & [-0.00003, +0.00024] \\
\midrule
Baby Products & R@20 & -0.00002 & [-0.00010, +0.00006] \\
Baby Products & nDCG@6 & 0.00000 & [-0.00003, +0.00002] \\
Baby Products & nDCG@10 & +0.00001 & [-0.00002, +0.00003] \\
Baby Products & nDCG@20 & 0.00000 & [-0.00002, +0.00002] \\
Baby Products & $U$ & -0.00002 & [-0.00007, +0.00003] \\
\midrule
Diginetica & R@20 & +0.00016 & [-0.00009, +0.00041] \\
Diginetica & nDCG@6 & +0.00005 & [-0.00006, +0.00015] \\
Diginetica & nDCG@10 & +0.00013 & [+0.00003, +0.00022] \\
Diginetica & nDCG@20 & +0.00007 & [0.00000, +0.00015] \\
Diginetica & $U$ & +0.00010 & [-0.00008, +0.00028] \\
\bottomrule
\end{tabular}
\end{table}

\begin{table}[H]
\centering
\scriptsize
\setlength{\tabcolsep}{3pt}
\caption{Primary linear-gate parameters on standardized OOF features, mean (sample SD) across seeds.}
\label{tab:primary-gate-coefficients}
\begin{tabular}{lrrrr}
\hline
Domain & Length & Frequency & Tail & Bias \\
\hline
Video Games & -0.07662 (0.04153) & -0.09143 (0.01712) & 0.01846 (0.01555) & -0.03097 (0.00398) \\
Baby Products & 0.01041 (0.03277) & -0.00834 (0.01360) & -0.04066 (0.02517) & -0.02598 (0.00267) \\
Diginetica & -0.05530 (0.01823) & -0.00348 (0.02062) & 0.04975 (0.05823) & -0.00651 (0.00603) \\
\hline
\end{tabular}
\end{table}


\section{Statistical Procedure}
For each method $m$, seed $s$, and query $q$, the retained per-query result
contains the top-20 target rank. Recall differences are exact paired values
$d_{qs}=\mathbf1[r_{m,qs}\leq K]-\mathbf1[r_{b,qs}\leq K]$.
We first average $d_{qs}$ over matched seeds, then sample query identifiers
with replacement. Each sampled query carries all three seed outcomes jointly,
preserving cross-seed pairing within query. Percentile endpoints from 20,000
replications form the 95\% interval. A sign-flip randomization test flips each
matched-seed per-query difference jointly. The evaluated HID Diginetica
representation contains expanded prefix rows but not their original session
identifier; the query is
therefore the smallest recoverable cluster and residual within-session
correlation is not modeled. These intervals describe recoverable-query
sampling conditional on the three fitted seeds; sample SD across those seeds
is reported separately.
The procedure is not a multiple-cutoff familywise test;
intervals at R@6/R@10/R@20 are descriptive matched analyses.

\section{Inference Runtime and Exact Replay}
Once the two top-120 expert lists are available, the gate is $O(3)$ and RRF
scores at most 240 distinct items before sorting that union. The added
allocation/fusion cost is therefore independent of catalogue size; full
catalogue scoring remains inside the frozen experts.
Features+gate and feature+gate+RRF columns in Table~\ref{tab:runtime} are
newly measured by exact NumPy inference. In a deterministic replay check,
reconstructed beta values differ from the original inference outputs by at
most \DynamicRuntimeMaxBetaError{}, and the top-20 rankings are reproduced
exactly. The total is explicitly a cross-run warm
component estimate: the newly measured dynamic post-processing is added to raw
CEARF retrieval and PASGR inference times from the previous audited warm-expert
run. It excludes every timing from the retired validation router.
\begin{table}[H]
\caption{Apple M2 Pro, 32\,GB, seed-42 frozen runs. Features+gate is included
inside Postproc.; Expert retrieval is CEARF memory retrieval plus PASGR
prediction from the prior warm-expert component run. Total = Expert retrieval
+ Postproc.; it is a cross-run component estimate, not a co-timed median.}
\label{tab:runtime}
\centering\scriptsize
\setlength{\tabcolsep}{2.5pt}
\begin{tabular}{lrrrrr}
\toprule
Domain & Queries & Feat.+gate ($\mu$s) & Postproc. ($\mu$s) &
Expert (ms) & Total (ms)\\
\midrule
Video Games & \DynamicRuntimeVideoQueries & \DynamicRuntimeVideoGateUs &
\DynamicRuntimeVideoPostUs & \DynamicRuntimeVideoExpertMs &
\DynamicRuntimeVideoTotalMs\\
Baby Products & \DynamicRuntimeBabyQueries & \DynamicRuntimeBabyGateUs &
\DynamicRuntimeBabyPostUs & \DynamicRuntimeBabyExpertMs &
\DynamicRuntimeBabyTotalMs\\
Diginetica & \DynamicRuntimeDigiQueries & \DynamicRuntimeDigiGateUs &
\DynamicRuntimeDigiPostUs & \DynamicRuntimeDigiExpertMs &
\DynamicRuntimeDigiTotalMs\\
\bottomrule
\end{tabular}
\end{table}

\section{Additional Implementation Details and Limitations}
All experiments use seeds 42, 123, and 456. The implementation follows the
source-disjoint OOF construction, allocator optimization, paired statistical
analysis, and inference procedures specified in Sections~3--15. All results
required to support the paper's claims are reported directly in the main
paper and this supplementary material.

The 5,000-query cap makes calibration affordable but may undersample rare
regimes. The single deterministic source holdout is out-of-fit, not full
$K$-fold cross-fitting. The primary gate uses only last-item popularity and
length context; richer uncertainty estimates may help on larger untouched
domains, but the current capacity audit warns against assuming that more
features are automatically safer. Finally, all main domains are
next-item benchmarks; long-horizon utility and online feedback remain open.

\enlargethispage{4\baselineskip}
\begin{thebibliography}{18}
\setlength{\itemsep}{-1pt}
\bibitem{cormack2009rrf} Cormack, G.V., Clarke, C.L.A., Buettcher, S.:
Reciprocal rank fusion outperforms Condorcet and individual rank learning
methods. In: SIGIR, pp. 758--759 (2009)
\bibitem{diginetica2016} DIGINETICA: CIKM Cup 2016 Track 2: Personalized
E-Commerce Search Challenge Data. CodaLab (2016)
\bibitem{garg2019stan} Garg, D., et al.: Sequence and time aware neighborhood
for session-based recommendations: STAN. In: SIGIR, pp. 1069--1072 (2019)
\bibitem{multichannelltr2026} Gaydhani, A., et al.: Unified learning-to-rank
for multi-channel retrieval in large-scale e-commerce search. CoRR
abs/2602.23530 (2026)
\bibitem{shortcut2026} Han, H., et al.: An embarrassingly simple graph
heuristic reveals shortcut-solvable benchmarks for sequential recommendation.
CoRR abs/2605.07125 (2026)
\bibitem{llm2rec2025} He, Y., et al.: LLM2Rec: Large language models are
powerful embedding models for sequential recommendation. In: KDD,
pp. 896--907 (2025)
\bibitem{hou2024amazonreviews} Hou, Y., et al.: Bridging language and items
for retrieval and recommendation. CoRR abs/2403.03952 (2024)
\bibitem{standards2026} Jannach, D., Chen, L.: Improving methodological
standards in recommender systems offline evaluation. ACM Trans. Recommender
Systems 4(3), 34e:1--34e:15 (2026)
\bibitem{li2017narm} Li, J., et al.: Neural attentive session-based
recommendation. In: CIKM, pp. 1419--1428 (2017)
\bibitem{li2025survey} Li, Z., et al.: Graph and sequential neural networks in
session-based recommendation: A survey. ACM Comput. Surv. 57(2), 1--37 (2025)
\bibitem{hm2rec2026} Liu, Y., et al.: Hyperbolic-enhanced mixture-of-experts
Mamba for sequential recommendation. In: AAAI 40, 15403--15411 (2026)
\bibitem{sigma2025} Liu, Z., et al.: SIGMA: Selective gated Mamba for
sequential recommendation. In: AAAI 39, 12264--12272 (2025)
\bibitem{ludewig2018vsknn} Ludewig, M., Jannach, D.: Evaluation of
session-based recommendation algorithms. UMUAI 28(4--5), 331--390 (2018)
\bibitem{wang2022e5} Wang, L., et al.: Text embeddings by weakly-supervised
contrastive pre-training. CoRR abs/2212.03533 (2022)
\bibitem{hidlongtail2026} Wang, X., et al.: Bid farewell to seesaw: Towards
accurate long-tail session-based recommendation via dual constraints of
hybrid intents. In: AAAI 40, 15895--15903 (2026)
\bibitem{wolpert1992stacking} Wolpert, D.H.: Stacked generalization. Neural
Networks 5(2), 241--259 (1992)
\bibitem{hidasi2016gru4rec} Hidasi, B., et al.: Session-based recommendations
with recurrent neural networks. In: ICLR (2016)
\bibitem{wu2019srgnn} Wu, S., et al.: Session-based recommendation with graph
neural networks. In: AAAI 33, 346--353 (2019)
\end{thebibliography}
\end{document}
